\chapter{Đặt vấn đề}

\section{Bối cảnh và động lực nghiên cứu}

Đi chung xe theo yêu cầu đã trở thành một trong những hướng được kỳ vọng nhiều nhất để giảm
áp lực giao thông đô thị. Ý tưởng rất trực tiếp: thay vì mỗi người một xe, hệ thống gom những
yêu cầu có lộ trình tương đồng vào cùng một phương tiện, nhờ đó giảm số xe lưu thông mà vẫn
phục vụ được cùng lượng nhu cầu. Công trình của Alonso-Mora và cộng sự \cite{alonsomora2017}
đã cho thấy tiềm năng của hướng này ở quy mô thật: với ba nghìn xe bốn chỗ, hệ thống của họ
phục vụ được khoảng $98\%$ nhu cầu taxi của Manhattan. Kết quả đó xác lập một chuẩn mực kỹ
thuật mà phần lớn các nghiên cứu sau này đều kế thừa, gồm đồ thị ba tầng Request--Trip--Vehicle
để sinh các khả năng ghép và một bài toán quy hoạch nguyên để chọn ra phương án tốt nhất.

Tuy nhiên, cách vận hành của các hệ thống này có một đặc điểm ít được nói tới. Chúng hoạt động
theo chu kỳ: cứ sau một khoảng thời gian ngắn, toàn bộ bài toán gán khách cho xe lại được giải
lại từ đầu với trạng thái mới. Mỗi lần giải lại, kế hoạch của một chiếc xe có thể thay đổi --
chèn thêm một điểm đón, đảo thứ tự hai điểm dừng, thậm chí chuyển một hành khách sang xe khác.
Điều đó là cần thiết để hệ thống tận dụng được mọi cơ hội ghép chuyến, và cũng chính là nguồn
gốc hiệu quả của nó.

Vấn đề nằm ở chỗ hành khách không nhìn thấy quá trình tối ưu ấy. Cái họ nhìn thấy là một con
số trên màn hình điện thoại: ``xe sẽ đón bạn lúc 9 giờ 10''. Khi hệ thống tối ưu lại và con số
đó trở thành 9 giờ 13, rồi 9 giờ 11, rồi 9 giờ 16, hành khách trải nghiệm một chuỗi lời hứa bị
rút lại chứ không phải một quá trình tối ưu. Và đây mới là điểm đáng chú ý: trong toàn bộ dòng
văn liệu hiện hành, chuỗi ấy \emph{hoàn toàn không được ghi nhận}. Các hệ thống này ràng buộc
hai tham số --- thời gian chờ tối đa $\Omega$ và độ trễ tối đa $\Delta$ --- và miễn là kế hoạch
còn nằm trong cửa sổ đó, nó được phép nhảy tùy ý, bao nhiêu lần cũng được. Một hành khách được
báo đúng một lần và một hành khách bị đổi giờ mười hai lần là hai đối tượng giống hệt nhau
trong mọi hàm mục tiêu của dòng bài này.

Sự thiếu vắng đó không phải chuyện nhỏ về mặt kinh tế. Một nhánh văn liệu khác, thuộc về phía
cầu, đã đo được rất rõ rằng hành khách trả giá cho độ tin cậy. Alonso-González và cộng sự
\cite{alonsogonzalez2020} ước lượng giá trị độ tin cậy cho dịch vụ đi chung theo yêu cầu;
Geržinič và cộng sự \cite{gerzinic2023} định giá riêng phần \emph{thời gian chờ bất ngờ} ---
tức phần chờ phát sinh ngoài dự kiến ban đầu --- ở mức cao hơn hẳn thời gian chờ thông thường.
Bansal và cộng sự \cite{bansal2019} cho thấy việc công bố ETA một cách trung thực làm tỉ lệ
chấp nhận chuyến tăng gần $10\%$. Nhưng chính các tác giả đó viết rõ trong bài rằng họ không
mô hình hóa chi phí phía cung: khung của họ giả định nhà điều hành \emph{luôn} có thể giao
khách đúng mốc đã hiển thị.

Nói cách khác, có cả một nhánh nghiên cứu chứng minh rằng giữ lời rất đáng tiền, và một nhánh
khác xây dựng những hệ thống điều phối hiệu quả, nhưng chưa ai trả lời câu hỏi nối hai nhánh
ấy lại:

\begin{khungy}
Giữ lời tốn bao nhiêu? Và liệu một lời hứa như thế có cưỡng chế được về mặt thuật toán hay
không, hay nó chỉ là một chỉ tiêu mong muốn?
\end{khungy}

\section{Hạn chế của các phương pháp hiện tại}

Để xác định chính xác chỗ đứng của khóa luận, chúng tôi đã khảo sát mười bốn trục văn liệu liên
quan. Từ đó rút ra bốn hạn chế mang tính hệ thống của dòng bài hiện hành, không phải những
thiếu sót ngẫu nhiên của một bài cụ thể.

\textbf{Thứ nhất, hành khách được mô hình hóa như một ràng buộc chứ không phải một đối tác.}
Cửa sổ $\Omega$ và $\Delta$ là tất cả những gì hệ thống nợ hành khách. Không tồn tại khái niệm
``điều đã nói với khách'' trong mô hình. Hệ quả trực tiếp là cửa sổ hợp lệ không hề bảo vệ lời
đã hứa: trong thực nghiệm của chúng tôi, một hành khách bị dời giờ trả tích lũy tới 573 giây
mà hệ thống vẫn tính là phục vụ thành công, vì mọi mốc vẫn nằm trong cửa sổ.

\textbf{Thứ hai, không có quy trách nguyên nhân.} Khi giờ hẹn của một hành khách xấu đi, không
ai tách phần do nhà điều hành quyết định lại khỏi phần do xe chạy chậm hoặc do đường tắc.
Fielbaum và Alonso-Mora \cite{fielbaum2020} là ngoại lệ đáng kể nhất --- họ phân loại nguyên
nhân thành ``do người dùng gây ra'' và ``do nhà điều hành gây ra'' --- nhưng đó là một phân
loại định tính bằng lời, và bản thân thước đo của họ thì loại tắc đường ra khỏi phạm vi ngay
từ đầu.

\textbf{Thứ ba, chỉ có thước đo kết quả, không có thước đo hành vi.} Bộ chỉ số chuẩn gồm tỉ lệ
phục vụ, thời gian chờ, quãng đường phương tiện và tỉ lệ ghép. Không chỉ số nào trả lời câu hỏi
``hệ thống này cư xử có giữ lời hay không''.

\textbf{Thứ tư, các khái niệm về ổn định lịch trình đã tồn tại ở ngành khác nhưng chưa được
chuyển sang lĩnh vực này.} Trong hoạch định sản xuất, hiện tượng ``MRP nervousness'' đã được
nghiên cứu từ thập niên 1980. Sridharan, Berry và Udayabhanu \cite{sridharan1987} đo độ ổn định
của lịch sản xuất chính dưới chân trời cuốn chiếu, và ``phương thuốc'' kinh điển của họ ---
đóng băng một phần lịch --- chính là cơ chế mà chúng tôi gọi là khóa lời hứa. Trong vận trù
giao hàng, bài toán Consistent VRP \cite{kovacs2014} có ràng buộc nhất quán giờ đến giữa các
ngày. Nhưng đó là bài toán tĩnh, nhiều ngày, cho hàng hóa; nó không nói gì về một lời hứa động
đưa ra cho một hành khách trong vòng vài phút.

\section{Chỗ đứng so với công trình gần nhất}

Trước khi trình bày đóng góp, cần nói rõ ranh giới với công trình gần nhất, vì phần chồng lấn
là đáng kể và việc giấu nó đi sẽ chỉ làm người đọc tự phát hiện ra ở chương sau.

Fielbaum và Alonso-Mora \cite{fielbaum2020} đã đặt đúng vấn đề nêu ở mục 1.1 và đã làm sáu việc
mà khóa luận này cũng làm: nêu khái niệm lời hứa công bố bị xê dịch, đo độ lớn dịch chuyển trên
từng hành khách, đếm số lần dự báo bị đổi, phân loại nguyên nhân thành do người dùng và do nhà
điều hành, dựng đường cong đánh đổi giữa mức kiểm soát và chất lượng dịch vụ, và thử cơ chế
ratchet trên giờ đón. Khóa luận coi sáu hạng mục ấy là \emph{nền}, và phần tương ứng trong
chương 4 là một cuộc \emph{tái lập độc lập} trên một thực thể khác --- có giá trị nhưng không
phải đóng góp mới.

Khác biệt nằm ở một điểm cấu trúc, và nó không phải chuyện mức độ mà là bản chất. Thước đo của
họ, công thức (2.5) ở chương sau, định nghĩa qua $t_w(r)$ --- thời gian chờ \emph{thực tế} ---
tức một đại lượng chỉ biết được sau khi hành khách đã được đón. Tại thời điểm hệ thống phải
chọn phương án nào, giá trị đó chưa tồn tại. Nghĩa là thước đo ấy, về mặt toán học, \emph{không
thể trở thành ràng buộc}: không ai chặn được ``thực tế trừ lần báo đầu không quá 30 giây'' khi
chưa biết thực tế. Nó chỉ tồn tại ở thì quá khứ.

\begin{khungy}
Công trình gần nhất \emph{mô tả} hiện tượng lời hứa bị xê dịch. Khóa luận này \emph{chặn} nó.
Để chặn được, đại lượng bị chặn phải định nghĩa trên những gì biết được tại thời điểm quyết
định --- và đó là lý do phép phân rã ba mốc ở chương 3 tồn tại.
\end{khungy}

Từ khác biệt cấu trúc ấy nảy sinh một câu hỏi mà trước đó không có lý do gì để đặt ra:
\emph{khi biến thước đo thành ràng buộc, nên tính sổ theo cái gì?} Chương 3 chứng minh rằng câu
trả lời quyết định hệ thống có vận hành nổi hay không, và rằng cách tự nhiên nhất để biến thước
đo của văn liệu thành ràng buộc lại chính là cách hỏng.

\section{Mục tiêu và đóng góp của khóa luận}

Từ bốn hạn chế ở mục 1.2 và khác biệt cấu trúc ở mục 1.3, khóa luận đặt mục tiêu xây dựng và
đánh giá một cơ chế cho phép hệ thống điều phối \emph{đưa ra lời hứa mà nó cưỡng chế được},
đồng thời đo chính xác cái giá của việc đó. Cụ thể gồm ba mục tiêu.

Mục tiêu thứ nhất là thiết kế cơ chế \textbf{lời hứa có biên được chứng nhận}: một sổ cái cam
kết riêng cho từng hành khách, cộng dồn theo thời gian và đối chiếu với một trần $\beta$, được
kiểm tra ngay tại thời điểm ra quyết định chứ không phải sau sự việc. Mục tiêu thứ hai là xây
dựng phép \textbf{phân rã ba mốc} để tách phần dịch chuyển do nhà điều hành gây ra khỏi phần
do thế giới vật lý gây ra, qua đó xác định chính xác cái gì nên bị tính vào sổ cái. Mục tiêu
thứ ba là đo bằng thực nghiệm có kiểm soát toàn bộ đường cong đánh đổi giữa mức độ chặt chẽ
của lời hứa và tỉ lệ phục vụ.

Các đóng góp chính của khóa luận gồm:

\begin{itemize}[leftmargin=1.2cm,itemsep=3pt,topsep=4pt]
\item \textbf{Ràng buộc cam kết cưỡng chế được.} Một sổ cái riêng cho từng hành khách, cộng dồn
theo chuỗi các lần công bố và đối chiếu với một trần $\beta$ ngay tại thời điểm quyết định.
Đây là điểm khác biệt cấu trúc đã nêu ở mục 1.3: đại lượng bị chặn định nghĩa hoàn toàn trên
những gì biết được trước khi chọn phương án, nên chặn được. Thực nghiệm xác nhận trần luôn được
tôn trọng ở mọi cấu hình --- $291 < 300$, $119 < 120$, $59 < 60$, $27{,}6 < 30$ giây.

\item \textbf{Kết quả cấu trúc về hai cơ sở tính sổ cái.} Chúng tôi chứng minh ba định lý cho
thấy cơ sở \emph{hành vi} luôn thỏa được, cơ sở \emph{kết quả} có thể khiến hệ thống không còn
quyết định hợp lệ nào, và cơ sở kết quả không bao giờ phục vụ được nhiều hơn cơ sở hành vi với
bất kỳ trần $\beta$ nào. Ba định lý này được kiểm chứng bằng thực nghiệm, bao gồm cả một thí
nghiệm bơm dữ liệu tắc đường đo thật của Manhattan, nơi cơ sở kết quả tê liệt trên 19 trên 20 ô
còn cơ sở hành vi không hỏng một ô nào.

\item \textbf{Phép phân rã ba mốc và quy trách định lượng.} Bằng cách chiếu lại chính kế hoạch
cũ từ vị trí thật hiện tại của xe, hệ thống tách được nguyên nhân của mỗi lần dời lời hứa theo
từng quyết định và từng hành khách. Đây là điều kiện kỹ thuật để hai đóng góp trên tồn tại, và
tự nó cho một con số: $89{,}06\%$ lượng dịch chuyển mà hành khách nhìn thấy trong chế độ không
cam kết là do bộ điều phối tự gây ra.

\item \textbf{Đường cong giá của cam kết.} Chín cấu hình cam kết khác nhau được đo trên cùng
một panel với cùng hạt giống, cho phép định lượng chi phí của từng mức chặt chẽ. Con số trung
tâm: một trần 30 giây cắt xáo trộn xuống còn $0{,}17\%$ mức ban đầu với giá $7{,}13$ điểm phần
trăm tỉ lệ phục vụ.

\item \textbf{Phân tích độ vững với cách đo.} Chín thước đo thuộc họ Sridharan (1988) được áp
dụng song song; thứ hạng các cấu hình không đổi dưới bất kỳ thước đo nào, trong khi độ lớn
chênh nhau tới hai bậc. Đi kèm là một cảnh báo phương pháp luận: tỉ lệ phục vụ nhiễu hơn nhiều
so với các thước đo ổn định, và một báo cáo chỉ dựa vào tỉ lệ phục vụ dễ đọc nhầm nhiễu thành
kết quả.

\item \textbf{Đối chiếu trực tiếp với Fielbaum và Alonso-Mora (2020).} Chúng tôi cài đặt lại
thước đo của họ trên dữ liệu của mình và chỉ ra rằng thước đo hai đầu của họ bỏ sót $49\%$ số
hành khách thực sự bị dời lời hứa dưới chế độ cam kết, đồng thời tái lập được kết luận định
tính trung tâm của họ trên một thực thể nhỏ hơn 185 lần.
\end{itemize}

Một điểm cần nói rõ ngay từ đầu, và sẽ được nhắc lại trong chương thực nghiệm: khóa luận này
\emph{không} tuyên bố vượt trội về tỉ lệ phục vụ. Mọi cấu hình có cam kết đều phục vụ ít hành
khách hơn cấu hình không cam kết. Đóng góp nằm ở chỗ khác --- ở việc biến một đại lượng trước
nay chỉ đo được sau sự việc thành một ràng buộc cưỡng chế được trước khi quyết định, và ở việc
chỉ ra rằng cách tính sổ cái quyết định hệ thống có vận hành nổi hay không.

\section{Phạm vi và giới hạn}

Khóa luận giới hạn phạm vi như sau. Dữ liệu là bộ chuyến đi taxi công khai của Manhattan
(nguồn TLC, giấy phép CC-BY) đã được chuẩn hóa sẵn trong bộ FleetPy \cite{fleetpy2022}, khung
thời gian 08:00--10:00 các ngày làm việc từ 14 đến 18 tháng 11 năm 2018. Quy mô mỗi ô thực
nghiệm là 108 yêu cầu, 8 xe sức chứa 4 chỗ, cửa sổ đón 600 giây, ràng buộc thời gian trên xe
bằng $1{,}5$ lần thời gian đi thẳng. Toàn bộ thực nghiệm dùng một hạt giống cố định nên hoàn
toàn tất định và tái lập được.

Quy mô này nhỏ hơn đáng kể so với các thực nghiệm của Alonso-Mora và cộng sự. Đó là lựa chọn
có chủ ý: mục tiêu của khóa luận là so sánh có kiểm soát giữa các cấu hình cam kết trên cùng
một đầu vào, chứ không phải chứng minh khả năng mở rộng. Hệ quả là các con số tuyệt đối trong
báo cáo này không so sánh trực tiếp được với con số trong các bài quy mô lớn; chỉ có các so
sánh nội bộ giữa các cấu hình mới có ý nghĩa. Chương 4 sẽ trở lại điểm này khi đối chiếu với
văn liệu.

\section{Cấu trúc báo cáo}

Phần còn lại của báo cáo được tổ chức theo đúng mạch lập luận đã nêu ở mục 1.3, tức đi từ chỗ
``thước đo không chặn được'' tới chỗ ``ràng buộc chặn được, và cái đó tốn bao nhiêu''.

Chương 2 trình bày cơ sở lý thuyết. Ngoài bài toán ghép chuyến và kiến trúc giải nó, chương này
dành một mục riêng để phân tích ba giới hạn của cách đo trong văn liệu, trong đó giới hạn thứ
hai --- thước đo chỉ tồn tại ở thì quá khứ --- là nền trực tiếp cho chương sau. Chương cũng vay
mượn hai khung khái niệm từ ngoài lĩnh vực: độ ổn định lịch trình trong hoạch định sản xuất, và
phân biệt nghĩa vụ phương tiện với nghĩa vụ kết quả trong luật hợp đồng.

Chương 3 là trọng tâm về phương pháp. Nó xây dựng phép phân rã ba mốc như lời giải cho vấn đề
``đại lượng bị chặn phải biết được tại thời điểm quyết định'', rồi từ đó dựng sổ cái cam kết,
và chứng minh ba định lý về hai cơ sở tính sổ. Chương kết thúc bằng kiến trúc cài đặt và cơ chế
chứng cứ niêm phong.

Chương 4 trình bày thực nghiệm, tổ chức theo khẳng định chứ không theo thứ tự các đợt chạy;
bảng lộ trình ở đầu chương nối mỗi khẳng định với mục chứng minh nó. Chương gồm cả phần đối
chiếu ranh giới với công trình gần nhất, phần ghi nhận các hướng đã thử và phải từ bỏ, và phần
liệt kê hạn chế.

Chương 5 tổng kết đóng góp, nêu rõ những hạn chế còn lại và đề xuất bốn hướng phát triển.
